% !TEX root = document.tex
\chapter{\label{chap:conclusion}Conclusion}

This thesis investigated whether higher-order pattern unification can be
combined with scope-graph-based name binding in the implementation of a
dependently typed language. To study this question, we implemented a small
dependently typed language in Statix and extended its type checker with higher-order inference. We subsequently implemented the
core of the same language in Elpi in order to compare the scope-graph
representation with Higher-Order Abstract Syntax.

With respect to RQ1, the Statix implementation shows that scope graphs can
provide the binding information required by higher-order pattern unification.
The variables available to a solution can be obtained from the
scope graph and used to construct well-scoped higher-order solutions.
Metavariables additionally require state that is visible independently of the
lexical scope in which a solution is discovered, which led to the separate
metavariable scope introduced in Chapter~5. The resulting implementation
supports the higher-order pattern cases considered in this thesis, although
the lack of constraint postponement restricts the set of unification problems
that can be handled.

With respect to RQ2, the comparison with Elpi shows that the two approaches
place complexity in different parts of the language implementation. HOAS
delegates lexical binding, capture-avoiding substitution, and higher-order
unification to the meta-language. Scope graphs require these operations to be
made more explicit, but provide very good support for non-lexical name
resolution, including modules, qualified references, and forward references.

Overall, this work shows that higher-order pattern unification is compatible
with a scope-graph-based representation of name binding. The main remaining
limitations arise not from scope graphs themselves, but from the additional
solver behavior required by a more complete dependent type elaborator.

\section{Future Work}
\label{sec:future-work}

There are several directions in which the implementation presented in this
thesis could be extended.

\paragraph{Constraint postponement and reactivation.}
The most immediate limitation of the current higher-order unification
procedure is the absence of constraint postponement. A constraint that is
outside the pattern fragment cannot necessarily be classified permanently as
unsolvable since later instantiations may simplify it into a pattern problem.
Supporting this behavior would require a mechanism for suspending
higher-order constraints and reactivating them when one of the metavariables
on which they depend is instantiated. Such a mechanism would considerably
reduce the number of constraints that the current implementation has to reject
eagerly.

\paragraph{Pruning and more complete higher-order pattern unification.}
Another possible extension is pruning. Sometimes a unification problem cannot
be solved because the proposed solution for a metavariable would refer to a
variable that was not in scope when that metavariable was introduced.
However, the offending variable may occur only as an argument to another metavariable. If that metavariable does not actually need to depend on the variable, its context can be reduced so that the offending dependency disappears. The original unification problem may then become solvable. Implementing this pruning step would allow the current solver to handle more higher-order constraints.

\paragraph{Closer integration with the Statix solver.}
The current implementation represents metavariable state explicitly through a
separate chain of scopes that is threaded through the type-checking rules.
This design was necessary because ordinary scope-graph bindings are immutable
and because the state of Statix unification variables cannot be inspected in
the way required by the implementation. An interesting direction would be to
investigate whether higher-order metavariables could be supported more
directly by the Statix solver. In particular, allowing a Statix unification variable to "see" a local context or exposing mechanisms for observing
the state of unification variables could remove part of the very explicit
state management currently required.

\section{Reproducibility}
The implementation artifacts developed as part of this thesis consist of an
Eclipse project containing the Statix specification of the dependently typed
language, together with a set of Elpi source files used for the alternative
HOAS-based implementation discussed in Chapter~6. They can be found here: \url{https://github.com/tdrAndrei/master_thesis/releases/tag/v1.0.0}.

The Statix artifact contains the language definition, type-checking rules,
scope-graph-based name resolution, metavariable handling, and higher-order
pattern unification described in Chapters~4 and~5. The project can be opened
in the Spoofax/Eclipse development environment and includes the examples and
tests used during development to validate the implementation. The Elpi
artifact contains the corresponding core language implementation used for the
comparison with Higher-Order Abstract Syntax.

\section{AI usage}
Generative AI tools were used during the preparation of this thesis, limited to tasks such as improving my own wording, formatting figures, LaTeX code and researching documentation or new papers to read.
They were not used as an authoritative source of scientific information.
AI was not used to develop the theory behind the metavariable scoping in Statix, or the Statix code. For the ELPI implementation of the language, it was used to find built-in functions that can apply multiple arguments to an ELPI metavariable.
In general, technical claims, definitions, algorithms, and comparisons included in this thesis were verified against the cited literature. 